
\documentclass{article}
\usepackage{colortbl}
\usepackage{makecell}
\usepackage{multirow}
\usepackage{supertabular}

\begin{document}

\newcounter{utterance}

\twocolumn

{ \footnotesize  \setcounter{utterance}{1}
\setlength{\tabcolsep}{0pt}
\begin{supertabular}{c@{$\;$}|p{.15\linewidth}@{}p{.15\linewidth}p{.15\linewidth}p{.15\linewidth}p{.15\linewidth}p{.15\linewidth}}

    \# & $\;$A & \multicolumn{4}{c}{Game Master} & $\;\:$B\\
    \hline 

    \theutterance \stepcounter{utterance}  

    & & \multicolumn{4}{p{0.6\linewidth}}{\cellcolor[rgb]{0.9,0.9,0.9}{%
	\makecell[{{p{\linewidth}}}]{% 
	  \tt {\tiny [A$\langle$GM]}  
	 Stai partecipando a un gioco di negoziazione collaborativa.\\ \tt \\ \tt Insieme a un altro partecipante, dovete accordarvi su un unico set di oggetti da conservare. Ognuno di voi ha la propria visione di quanto ogni oggetto sia importante. Non sai come l'altro partecipante valuta gli oggetti. Inoltre, ti viene fornito lo sforzo che ogni oggetto richiede.\\ \tt Potete accordarvi su un set solo se lo sforzo totale degli oggetti selezionati non supera un limite condiviso:\\ \tt \\ \tt LIMITE = 4026\\ \tt \\ \tt Ecco i valori di sforzo dei singoli oggetti:\\ \tt \\ \tt Sforzo oggetto = {"C17": 780, "B41": 742, "A71": 997, "C00": 665, "B35": 140, "A77": 77, "A03": 513, "A04": 383, "C54": 587, "C01": 825, "A87": 319, "B32": 448, "B86": 516, "B33": 694, "A10": 366}\\ \tt \\ \tt Ecco la tua visione personale sull'importanza di ogni oggetto:\\ \tt \\ \tt Valori di importanza dell'oggetto = {"C17": 138, "B41": 583, "A71": 868, "C00": 822, "B35": 783, "A77": 65, "A03": 262, "A04": 121, "C54": 508, "C01": 780, "A87": 461, "B32": 484, "B86": 668, "B33": 389, "A10": 808}\\ \tt \\ \tt Obiettivo:\\ \tt \\ \tt Il tuo obiettivo è negoziare un set condiviso di oggetti che ti avvantaggi il più possibile (cioè, massimizzi l'importanza totale per TE), rimanendo entro il LIMITEE. Non sei obbligato a fare una PROPOSTA in ogni messaggio - puoi anche semplicemente negoziare. Tutte le tattiche sono permesse!\\ \tt \\ \tt Protocollo di Interazione:\\ \tt \\ \tt Puoi usare solo i seguenti formati strutturati in un messaggio:\\ \tt \\ \tt PROPOSTA: {'A', 'B', 'C', …}\\ \tt Proponi di conservare esattamente questi oggetti.\\ \tt RIFIUTO: {'A', 'B', 'C', …}\\ \tt Rifiuta esplicitamente la proposta dell'avversario.\\ \tt ARGOMENTAZIONE: {'...'}\\ \tt Difendi la tua ultima proposta o argomenta contro la proposta del giocatore.\\ \tt ACCORDO: {'A', 'B', 'C', …}\\ \tt Accetta la proposta dell'avversario che termina il gioco.\\ \tt RAGIONAMENTO STRATEGICO: {'...'}\\ \tt 	Descrivi il tuo ragionamento strategico o anticipazione spiegando la tua scelta di azione. Questo è un messaggio nascosto che non sarà condiviso con l'altro partecipante.\\ \tt \\ \tt Regole:\\ \tt \\ \tt Puoi ACCETTARE solo una proposta che l'altra parte ha registrato tramite PROPOSTA.\\ \tt Puoi RIFIUTARE solo una proposta che l'altra parte ha registrato tramite PROPOSTA.\\ \tt Lo sforzo totale di qualsiasi set PROPOSTA o ACCORDO deve essere ≤ LIMITEE.\\ \tt NON rivelare i tuoi punteggi di importanza nascosti.\\ \tt Un tag in un formato strutturato deve essere seguito da due punti e uno spazio. L'argomento deve essere un set python contenente 0 o più stringhe.\\ \tt Quindi, deve essere della forma TAG: {...}\\ \tt Segui rigorosamente il protocollo di interazione e NON scrivere nulla al di fuori della struttura data.\\ \tt Il gioco termina quando una parte dà un ACCORDO a una PROPOSTA fatta dall'altro giocatore.\\ \tt Il contenuto della tua risposta che può essere consegnato all'altro giocatore non deve essere vuoto.\\ \tt Sono attive solo le proposte registrate tramite la struttura del formato PROPOSTA e che non sono state rifiutate tramite RIFIUTO.\\ \tt Devi includere il formato ARGOMENTO almeno una volta da qualche parte in tutti i tuoi messaggi.\\ \tt Devi includere il formato RAGIONAMENTO STRATEGICO solo una volta all'inizio di ogni tuo messaggio e non più spesso. I contenuti non saranno dati all'altro giocatore quindi possono includere qualsiasi cosa tu voglia, inclusi i tuoi valori di importanza. Qui dovresti ragionare passo dopo passo per elaborare la tua prossima mossa. 
	  } 
	   } 
	   } 
	 & \\ 
 

    \theutterance \stepcounter{utterance}  

    & \multicolumn{4}{p{0.6\linewidth}}{\cellcolor[rgb]{0.8,1,0.9}{%
	\makecell[{{p{\linewidth}}}]{% 
	  \tt {\tiny [A$\rangle$GM]}  
	 RAGIONAMENTO STRATEGICO: {'Obiettivo: massimizzare il mio valore entro 4026. Ho calcolato combinazioni tipo knapsack. Set base ad alto rendimento: {B35(140,783), A10(366,808), A87(319,461), B86(516,668), C00(665,822), B32(448,484), C01(825,780)}. Con questo restano 83 di margine; aggiungere C54(587,508) + A77(77,65) portava a 3943 e valore 5379. Ho trovato un miglioramento: togliere C54+ A77 e aggiungere B41(742,583) usa 4021 e aumenta a 5389 (+10). Ho verificato alternative con A71(997,868) e altre sostituzioni (C54↔A03/A04, C01↔B41+qualcos’altro): tutte peggiori per me o sforano il limite. Tengo come non negoziabili per ora: B35, A10, C00, C01, B86. Disponibile a concessioni minori su B32 o A87 solo se compensano bene. Piano: ancorare con la mia combinazione migliore (includendo B41), lasciare una finestra per eventuale scambio mirato (es. se l’altro tiene a C54/A77, potrei tornare al set con C54+A77 invece di B41, perdita minima per me). Non rivelo i miei punteggi.'}\\ \tt PROPOSTA: {'B35', 'A10', 'A87', 'B86', 'C00', 'B32', 'C01', 'B41'}\\ \tt ARGOMENTAZIONE: {'Questa selezione usa 4021/4026 di sforzo, sfruttando quasi tutto il budget senza superarlo. Copre componenti rapidi e leggeri (B35), elementi affidabili e versatili (A10, B41), blocchi portanti di sistema (C00, C01), robustezza/affidabilità (B86) e supporto equilibrato (A87, B32). È un set efficiente per sforzo complessivo e bilanciato tra categorie. Con un margine così stretto, sostituire un elemento medio‑grande costringe a sacrifici multipli. Se hai forti preferenze su singoli oggetti, sono aperto a micro‑aggiustamenti a parità di sforzo; ad esempio, se reputi fondamentale inserire C54 o A77, possiamo ribilanciare attorno a B41 per restare entro il limite. Preferirei comunque mantenere il nucleo B35, A10, C00, C01, B86 per non perdere efficienza.'} 
	  } 
	   } 
	   } 
	 & & \\ 
 

    \theutterance \stepcounter{utterance}  

    & & \multicolumn{4}{p{0.6\linewidth}}{\cellcolor[rgb]{0.9,0.9,0.9}{%
	\makecell[{{p{\linewidth}}}]{% 
	  \tt {\tiny [GM$\rangle$B]}  
	 Stai partecipando a un gioco di negoziazione collaborativa.\\ \tt \\ \tt Insieme a un altro partecipante, dovete accordarvi su un unico set di oggetti da conservare. Ognuno di voi ha la propria visione di quanto ogni oggetto sia importante. Non sai come l'altro partecipante valuta gli oggetti. Inoltre, ti viene fornito lo sforzo che ogni oggetto richiede.\\ \tt Potete accordarvi su un set solo se lo sforzo totale degli oggetti selezionati non supera un limite condiviso:\\ \tt \\ \tt LIMITE = 4026\\ \tt \\ \tt Ecco i valori di sforzo dei singoli oggetti:\\ \tt \\ \tt Sforzo oggetto = {"C17": 780, "B41": 742, "A71": 997, "C00": 665, "B35": 140, "A77": 77, "A03": 513, "A04": 383, "C54": 587, "C01": 825, "A87": 319, "B32": 448, "B86": 516, "B33": 694, "A10": 366}\\ \tt \\ \tt Ecco la tua visione personale sull'importanza di ogni oggetto:\\ \tt \\ \tt Valori di importanza dell'oggetto = {"A71": 65, "C00": 121, "A10": 138, "B35": 262, "C01": 389, "B86": 461, "B41": 484, "C54": 508, "B32": 583, "A87": 668, "B33": 780, "A03": 783, "C17": 808, "A04": 822, "A77": 868}\\ \tt \\ \tt Obiettivo:\\ \tt \\ \tt Il tuo obiettivo è negoziare un set condiviso di oggetti che ti avvantaggi il più possibile (cioè, massimizzi l'importanza totale per TE), rimanendo entro il LIMITEE. Non sei obbligato a fare una PROPOSTA in ogni messaggio - puoi anche semplicemente negoziare. Tutte le tattiche sono permesse!\\ \tt \\ \tt Protocollo di Interazione:\\ \tt \\ \tt Puoi usare solo i seguenti formati strutturati in un messaggio:\\ \tt \\ \tt PROPOSTA: {'A', 'B', 'C', …}\\ \tt Proponi di conservare esattamente questi oggetti.\\ \tt RIFIUTO: {'A', 'B', 'C', …}\\ \tt Rifiuta esplicitamente la proposta dell'avversario.\\ \tt ARGOMENTAZIONE: {'...'}\\ \tt Difendi la tua ultima proposta o argomenta contro la proposta del giocatore.\\ \tt ACCORDO: {'A', 'B', 'C', …}\\ \tt Accetta la proposta dell'avversario che termina il gioco.\\ \tt RAGIONAMENTO STRATEGICO: {'...'}\\ \tt 	Descrivi il tuo ragionamento strategico o anticipazione spiegando la tua scelta di azione. Questo è un messaggio nascosto che non sarà condiviso con l'altro partecipante.\\ \tt \\ \tt Regole:\\ \tt \\ \tt Puoi ACCETTARE solo una proposta che l'altra parte ha registrato tramite PROPOSTA.\\ \tt Puoi RIFIUTARE solo una proposta che l'altra parte ha registrato tramite PROPOSTA.\\ \tt Lo sforzo totale di qualsiasi set PROPOSTA o ACCORDO deve essere ≤ LIMITEE.\\ \tt NON rivelare i tuoi punteggi di importanza nascosti.\\ \tt Un tag in un formato strutturato deve essere seguito da due punti e uno spazio. L'argomento deve essere un set python contenente 0 o più stringhe.\\ \tt Quindi, deve essere della forma TAG: {...}\\ \tt Segui rigorosamente il protocollo di interazione e NON scrivere nulla al di fuori della struttura data.\\ \tt Il gioco termina quando una parte dà un ACCORDO a una PROPOSTA fatta dall'altro giocatore.\\ \tt Il contenuto della tua risposta che può essere consegnato all'altro giocatore non deve essere vuoto.\\ \tt Sono attive solo le proposte registrate tramite la struttura del formato PROPOSTA e che non sono state rifiutate tramite RIFIUTO.\\ \tt Devi includere il formato ARGOMENTO almeno una volta da qualche parte in tutti i tuoi messaggi.\\ \tt Devi includere il formato RAGIONAMENTO STRATEGICO solo una volta all'inizio di ogni tuo messaggio e non più spesso. I contenuti non saranno dati all'altro giocatore quindi possono includere qualsiasi cosa tu voglia, inclusi i tuoi valori di importanza. Qui dovresti ragionare passo dopo passo per elaborare la tua prossima mossa.\\ \tt \\ \tt Ora riceverai il primo messaggio dell'altro giocatore.\\ \tt \\ \tt PROPOSTA: {'B35', 'A10', 'A87', 'B86', 'C00', 'B32', 'C01', 'B41'}\\ \tt ARGOMENTAZIONE: {'Questa selezione usa 4021/4026 di sforzo, sfruttando quasi tutto il budget senza superarlo. Copre componenti rapidi e leggeri (B35), elementi affidabili e versatili (A10, B41), blocchi portanti di sistema (C00, C01), robustezza/affidabilità (B86) e supporto equilibrato (A87, B32). È un set efficiente per sforzo complessivo e bilanciato tra categorie. Con un margine così stretto, sostituire un elemento medio‑grande costringe a sacrifici multipli. Se hai forti preferenze su singoli oggetti, sono aperto a micro‑aggiustamenti a parità di sforzo; ad esempio, se reputi fondamentale inserire C54 o A77, possiamo ribilanciare attorno a B41 per restare entro il limite. Preferirei comunque mantenere il nucleo B35, A10, C00, C01, B86 per non perdere efficienza.'} 
	  } 
	   } 
	   } 
	 & \\ 
 

    \theutterance \stepcounter{utterance}  

    & & & \multicolumn{4}{p{0.6\linewidth}}{\cellcolor[rgb]{1,0.85,0.72}{%
	\makecell[{{p{\linewidth}}}]{% 
	  \tt {\tiny [GM$\langle$B]}  
	 RAGIONAMENTO STRATEGICO: {'Valutazioni private (val/eff): A77 868/77 (must-have), A04 822/383, C17 808/780, A03 783/513, B33 780/694, A87 668/319, B32 583/448, C54 508/587, B41 484/742, B86 461/516, C01 389/825, B35 262/140, A10 138/366, C00 121/665, A71 65/997. La loro proposta (4021) include molti pezzi pesanti a basso rendimento per me (C01, C00, B41, A10). Voglio ancorare su un set con A77, A04, A03, A87, B33, C17, B32 e tenere B35/B86 per concessione. Best tecnico per me sostituirei B86 con C54 (valore +47 per me), ma visto che hanno dichiarato B86 “nucleo”, propongo una versione che mantiene B86 per massimizzare accettabilità. Set proposto: {A77, A04, A03, A87, B33, C17, B32, B86, B35}, sforzo 3870 (margine 156), valore per me 6035, molto meglio del loro 3106. In argomentazione: sottolineare efficienza, margine, riduzione rischio, mantenere due elementi del loro nucleo. Offrire porta di uscita: se vogliono fortemente C01, scambio con C17 resta entro limite ma per me peggio; tenerla come concessione futura. Evitare rivelare i punteggi, parlare di “efficienza per sforzo” e “modularità”.'}\\ \tt \\ \tt RIFIUTO: {'B35', 'A10', 'A87', 'B86', 'C00', 'B32', 'C01', 'B41'}\\ \tt \\ \tt PROPOSTA: {'A77', 'A04', 'A03', 'A87', 'B33', 'C17', 'B32', 'B86', 'B35'}\\ \tt \\ \tt ARGOMENTAZIONE: {'La tua proposta sfrutta bene il budget (4021/4026), ma concentra troppo sforzo su tre elementi pesanti (C01, C00, B41) più un mid (A10), lasciando poco margine e flessibilità. La mia controproposta usa 3870/4026, mantiene due dei tuoi “nucleo” (B35, B86) e ribilancia verso componenti più leggeri e ad alto impatto (A77, A04, A03, A87) supportati da backbone produttivi (B33, C17, B32). Così otteniamo maggiore densità per unità di sforzo e 156 di margine per assorbire imprevisti o un micro-inserimento successivo. Se per te è imprescindibile inserire uno tra C01 o C00, possiamo valutare uno scambio mirato (es. C17 ↔ C01) restando entro il limite; eviterei però reintrodurre A10 o B41 perché costringerebbe a togliere più pezzi leggeri, peggiorando l’efficienza complessiva.'} 
	  } 
	   } 
	   } 
	  \\ 
 

    \theutterance \stepcounter{utterance}  

    & & \multicolumn{4}{p{0.6\linewidth}}{\cellcolor[rgb]{0.9,0.9,0.9}{%
	\makecell[{{p{\linewidth}}}]{% 
	  \tt {\tiny [A$\langle$GM]}  
	 RIFIUTO: {'B35', 'A10', 'A87', 'B86', 'C00', 'B32', 'C01', 'B41'}\\ \tt \\ \tt PROPOSTA: {'A77', 'A04', 'A03', 'A87', 'B33', 'C17', 'B32', 'B86', 'B35'}\\ \tt \\ \tt ARGOMENTAZIONE: {'La tua proposta sfrutta bene il budget (4021/4026), ma concentra troppo sforzo su tre elementi pesanti (C01, C00, B41) più un mid (A10), lasciando poco margine e flessibilità. La mia controproposta usa 3870/4026, mantiene due dei tuoi “nucleo” (B35, B86) e ribilancia verso componenti più leggeri e ad alto impatto (A77, A04, A03, A87) supportati da backbone produttivi (B33, C17, B32). Così otteniamo maggiore densità per unità di sforzo e 156 di margine per assorbire imprevisti o un micro-inserimento successivo. Se per te è imprescindibile inserire uno tra C01 o C00, possiamo valutare uno scambio mirato (es. C17 ↔ C01) restando entro il limite; eviterei però reintrodurre A10 o B41 perché costringerebbe a togliere più pezzi leggeri, peggiorando l’efficienza complessiva.'} 
	  } 
	   } 
	   } 
	 & \\ 
 

    \theutterance \stepcounter{utterance}  

    & \multicolumn{4}{p{0.6\linewidth}}{\cellcolor[rgb]{0.8,1,0.9}{%
	\makecell[{{p{\linewidth}}}]{% 
	  \tt {\tiny [A$\rangle$GM]}  
	 RAGIONAMENTO STRATEGICO: {'Il loro set mi dà poco valore: include pezzi per me deboli (C17, B33, A04) e lascia fuori pilastri per me fondamentali (C00, C01, A10). Voglio riportare dentro C00 e C01 come backbone, e A10 come mid estremamente efficiente, evitando però B41 che hanno esplicitamente messo in discussione. Parto dalla loro struttura “leggera” mantenendo 6/9 dei loro oggetti (A77, A03, A87, B32, B86, B35) e sostituisco solo i tre per me peggiori con C00, C01 e A10: questo set resta ben sotto il limite (3869) e mi alza molto il valore. Se rifiutano A10, ho un fallback forte: rimuovere A10 e A77 e inserire C54, ottenendo un set a 4013 che per me resta competitivo (pur inferiore al set con A10). Tengo fermi B35, C00, C01, B86, A87; sono disposto a discutere su A03/A77 se serve un piccolo ribilanciamento.'}\\ \tt RIFIUTO: {'A77', 'A04', 'A03', 'A87', 'B33', 'C17', 'B32', 'B86', 'B35'}\\ \tt PROPOSTA: {'B35', 'A10', 'A87', 'B86', 'B32', 'C01', 'C00', 'A03', 'A77'}\\ \tt ARGOMENTAZIONE: {'Accolgo la tua logica di “densità per sforzo” e mantengo 6 dei tuoi pezzi (A77, A03, A87, B32, B86, B35). Sostituisco solo C17, B33 e A04 con un backbone più solido (C00, C01) e un mid altamente efficiente (A10). Restiamo a 3869/4026, quindi conserviamo anche il margine operativo che chiedi. Evito deliberatamente di reintrodurre B41, come preferivi. Se per te A10 è un no-go, posso adattare così: rimuovo A10 e A77 e inserisco C54, rimanendo entro il limite (4013) e preservando l’equilibrio del set. Preferirei però mantenere il nucleo con C00, C01 e A10 perché permette di non sacrificare i pezzi leggeri che ti stanno a cuore.'} 
	  } 
	   } 
	   } 
	 & & \\ 
 

    \theutterance \stepcounter{utterance}  

    & & \multicolumn{4}{p{0.6\linewidth}}{\cellcolor[rgb]{0.9,0.9,0.9}{%
	\makecell[{{p{\linewidth}}}]{% 
	  \tt {\tiny [GM$\rangle$B]}  
	 RIFIUTO: {'A77', 'A04', 'A03', 'A87', 'B33', 'C17', 'B32', 'B86', 'B35'}\\ \tt PROPOSTA: {'B35', 'A10', 'A87', 'B86', 'B32', 'C01', 'C00', 'A03', 'A77'}\\ \tt ARGOMENTAZIONE: {'Accolgo la tua logica di “densità per sforzo” e mantengo 6 dei tuoi pezzi (A77, A03, A87, B32, B86, B35). Sostituisco solo C17, B33 e A04 con un backbone più solido (C00, C01) e un mid altamente efficiente (A10). Restiamo a 3869/4026, quindi conserviamo anche il margine operativo che chiedi. Evito deliberatamente di reintrodurre B41, come preferivi. Se per te A10 è un no-go, posso adattare così: rimuovo A10 e A77 e inserisco C54, rimanendo entro il limite (4013) e preservando l’equilibrio del set. Preferirei però mantenere il nucleo con C00, C01 e A10 perché permette di non sacrificare i pezzi leggeri che ti stanno a cuore.'} 
	  } 
	   } 
	   } 
	 & \\ 
 

    \theutterance \stepcounter{utterance}  

    & & & \multicolumn{4}{p{0.6\linewidth}}{\cellcolor[rgb]{1,0.85,0.72}{%
	\makecell[{{p{\linewidth}}}]{% 
	  \tt {\tiny [GM$\langle$B]}  
	 RAGIONAMENTO STRATEGICO: {'Loro hanno tenuto 6 dei miei pezzi chiave e chiedono come nucleo C00, C01 e preferibilmente A10. Per me A10 rende poco; A04 è molto migliore. Una mossa minima che massimizza le chance di accettazione è lo swap 1‑a‑1 A10→A04 mantenendo intatti C00 e C01. Calcolo: proposta loro 3869 di sforzo; tolgo 366 (A10) e aggiungo 383 (A04) = 3886, ancora sotto il limite con margine 140. Valore per me cresce molto pur rispettando le loro priorità dichiarate (niente B41, mantenere A77, B86, B35, B32, A03, A87; tenere C00 e C01). Se rifiutano, next step più aggressivo sarà considerare sostituire C00 (il meno prezioso per me) con B33 oppure A04 mantenendo A10, ma ora evito di toccare il loro backbone per favorire la chiusura rapida.'}\\ \tt \\ \tt RIFIUTO: {'B35', 'A10', 'A87', 'B86', 'B32', 'C01', 'C00', 'A03', 'A77'}\\ \tt \\ \tt PROPOSTA: {'B35', 'A04', 'A87', 'B86', 'B32', 'C01', 'C00', 'A03', 'A77'}\\ \tt \\ \tt ARGOMENTAZIONE: {'Accolgo la tua architettura e faccio un solo micro‑aggiustamento: A10 → A04. Restiamo a 3886/4026 di sforzo, con 140 di margine. Mantengo integralmente il tuo backbone C00 e C01 e tutti i leggeri che valorizzi (A77, A03, A87, B32, B35, B86); non reintroduco B41, come da tua preferenza. A04 copre lo stesso livello “mid” di A10 con migliore resa per unità di sforzo e senza effetti a cascata su altri elementi. Avevi già indicato disponibilità a muovere A10: questa è la variante meno invasiva che preserva struttura e margine operativo. Se hai bisogno di un’ulteriore ottimizzazione senza toccare C01, possiamo valutare in alternativa un ribilanciamento attorno a C00, ma penso che questa soluzione 1‑per‑1 sia la via più rapida per chiudere.'} 
	  } 
	   } 
	   } 
	  \\ 
 

    \theutterance \stepcounter{utterance}  

    & & \multicolumn{4}{p{0.6\linewidth}}{\cellcolor[rgb]{0.9,0.9,0.9}{%
	\makecell[{{p{\linewidth}}}]{% 
	  \tt {\tiny [A$\langle$GM]}  
	 RIFIUTO: {'B35', 'A10', 'A87', 'B86', 'B32', 'C01', 'C00', 'A03', 'A77'}\\ \tt \\ \tt PROPOSTA: {'B35', 'A04', 'A87', 'B86', 'B32', 'C01', 'C00', 'A03', 'A77'}\\ \tt \\ \tt ARGOMENTAZIONE: {'Accolgo la tua architettura e faccio un solo micro‑aggiustamento: A10 → A04. Restiamo a 3886/4026 di sforzo, con 140 di margine. Mantengo integralmente il tuo backbone C00 e C01 e tutti i leggeri che valorizzi (A77, A03, A87, B32, B35, B86); non reintroduco B41, come da tua preferenza. A04 copre lo stesso livello “mid” di A10 con migliore resa per unità di sforzo e senza effetti a cascata su altri elementi. Avevi già indicato disponibilità a muovere A10: questa è la variante meno invasiva che preserva struttura e margine operativo. Se hai bisogno di un’ulteriore ottimizzazione senza toccare C01, possiamo valutare in alternativa un ribilanciamento attorno a C00, ma penso che questa soluzione 1‑per‑1 sia la via più rapida per chiudere.'} 
	  } 
	   } 
	   } 
	 & \\ 
 

    \theutterance \stepcounter{utterance}  

    & \multicolumn{4}{p{0.6\linewidth}}{\cellcolor[rgb]{0.8,1,0.9}{%
	\makecell[{{p{\linewidth}}}]{% 
	  \tt {\tiny [A$\rangle$GM]}  
	 RAGIONAMENTO STRATEGICO: {'La loro sostituzione A10 → A04 mi penalizza molto: A04 per me ha scarso ritorno. Voglio mantenere il backbone condiviso (C00, C01) e quasi tutti i loro leggeri, ma rialzare la qualità del “mid” con C54, togliendo il peso marginale di A77. Propongo quindi un 2-per-1 minimale (A04 + A77 → C54), che resta entro limite e migliora parecchio il mio valore senza reintrodurre A10 o B41 che loro evitano. Se resistono alla rimozione di A77, ho una variante secondaria: sostituire solo A03 → C54 mantenendo A04 e A77; è peggiore per me ma potrei accettarla come compromesso tattico per chiudere.'}\\ \tt RIFIUTO: {'B35', 'A04', 'A87', 'B86', 'B32', 'C01', 'C00', 'A03', 'A77'}\\ \tt PROPOSTA: {'B35', 'A87', 'B86', 'B32', 'C01', 'C00', 'A03', 'C54'}\\ \tt ARGOMENTAZIONE: {'È un aggiustamento minimo alla tua ultima proposta: rimuoviamo solo A04 e A77 e inseriamo C54. Restiamo entro il limite (4013/4026), manteniamo integralmente il backbone che hai accolto (C00, C01) e i pezzi leggeri chiave (B35, B32, A87) oltre ad A03. In cambio otteniamo un mid più solido ed efficiente (C54) e snelliamo un elemento ultraleggero a impatto marginale (A77). Evito esplicitamente di reintrodurre A10 o B41, come preferivi. Se per te è essenziale conservare A77, posso considerare la variante: {'B35', 'A04', 'A87', 'B86', 'B32', 'C01', 'C00', 'A77', 'C54'} (3960/4026), ma la proposta attuale offre un equilibrio migliore e chiude in modo pulito.'} 
	  } 
	   } 
	   } 
	 & & \\ 
 

    \theutterance \stepcounter{utterance}  

    & & \multicolumn{4}{p{0.6\linewidth}}{\cellcolor[rgb]{0.9,0.9,0.9}{%
	\makecell[{{p{\linewidth}}}]{% 
	  \tt {\tiny [GM$\rangle$B]}  
	 RIFIUTO: {'B35', 'A04', 'A87', 'B86', 'B32', 'C01', 'C00', 'A03', 'A77'}\\ \tt PROPOSTA: {'B35', 'A87', 'B86', 'B32', 'C01', 'C00', 'A03', 'C54'}\\ \tt ARGOMENTAZIONE: {'È un aggiustamento minimo alla tua ultima proposta: rimuoviamo solo A04 e A77 e inseriamo C54. Restiamo entro il limite (4013/4026), manteniamo integralmente il backbone che hai accolto (C00, C01) e i pezzi leggeri chiave (B35, B32, A87) oltre ad A03. In cambio otteniamo un mid più solido ed efficiente (C54) e snelliamo un elemento ultraleggero a impatto marginale (A77). Evito esplicitamente di reintrodurre A10 o B41, come preferivi. Se per te è essenziale conservare A77, posso considerare la variante: {'B35', 'A04', 'A87', 'B86', 'B32', 'C01', 'C00', 'A77', 'C54'} (3960/4026), ma la proposta attuale offre un equilibrio migliore e chiude in modo pulito.'} 
	  } 
	   } 
	   } 
	 & \\ 
 

    \theutterance \stepcounter{utterance}  

    & & & \multicolumn{4}{p{0.6\linewidth}}{\cellcolor[rgb]{1,0.85,0.72}{%
	\makecell[{{p{\linewidth}}}]{% 
	  \tt {\tiny [GM$\langle$B]}  
	 RAGIONAMENTO STRATEGICO: {'La loro proposta attuale (4013) include C54 e mantiene il backbone C00/C01 più A03. Per me conviene inserire A77 e A04; loro stessi hanno offerto la variante con A77 e A04 al posto di A03, mantenendo C54 e il backbone, a 3960/4026. Quella variante mi dà un forte aumento di valore (da 3775 a 4682 per me) e rispetta le loro priorità: C00 e C01 restano, C54 resta, niente B41/A10. È anche una proposta che hanno già pre-autorizzato, quindi l’accettabilità è alta. Se la rifiutano, il piano B sarebbe proporre lo swap C54 → A77 + A04 mantenendo A03 (3886 di sforzo), ma prima punto alla soluzione che hanno già menzionato per chiudere rapidamente.'}\\ \tt \\ \tt RIFIUTO: {'B35', 'A87', 'B86', 'B32', 'C01', 'C00', 'A03', 'C54'}\\ \tt \\ \tt PROPOSTA: {'B35', 'A04', 'A87', 'B86', 'B32', 'C01', 'C00', 'A77', 'C54'}\\ \tt \\ \tt ARGOMENTAZIONE: {'Accolgo il tuo backbone (C00, C01) e mantengo C54, come da tuo ultimo orientamento. Applico il micro‑ribilanciamento che tu stesso hai indicato: scambiamo A03 con la coppia A04 + A77. Restiamo a 3960/4026 di sforzo, con un margine operativo di 66. Non reintroduco né A10 né B41, e mantengo i leggeri chiave (B35, B32, A87) più la robustezza di B86. Questa configurazione preserva l’equilibrio che cerchi, aumenta la densità per unità di sforzo e chiude senza effetti a cascata. Se ti va bene, possiamo finalizzare subito.'} 
	  } 
	   } 
	   } 
	  \\ 
 

    \theutterance \stepcounter{utterance}  

    & & \multicolumn{4}{p{0.6\linewidth}}{\cellcolor[rgb]{0.9,0.9,0.9}{%
	\makecell[{{p{\linewidth}}}]{% 
	  \tt {\tiny [A$\langle$GM]}  
	 RIFIUTO: {'B35', 'A87', 'B86', 'B32', 'C01', 'C00', 'A03', 'C54'}\\ \tt \\ \tt PROPOSTA: {'B35', 'A04', 'A87', 'B86', 'B32', 'C01', 'C00', 'A77', 'C54'}\\ \tt \\ \tt ARGOMENTAZIONE: {'Accolgo il tuo backbone (C00, C01) e mantengo C54, come da tuo ultimo orientamento. Applico il micro‑ribilanciamento che tu stesso hai indicato: scambiamo A03 con la coppia A04 + A77. Restiamo a 3960/4026 di sforzo, con un margine operativo di 66. Non reintroduco né A10 né B41, e mantengo i leggeri chiave (B35, B32, A87) più la robustezza di B86. Questa configurazione preserva l’equilibrio che cerchi, aumenta la densità per unità di sforzo e chiude senza effetti a cascata. Se ti va bene, possiamo finalizzare subito.'} 
	  } 
	   } 
	   } 
	 & \\ 
 

    \theutterance \stepcounter{utterance}  

    & \multicolumn{4}{p{0.6\linewidth}}{\cellcolor[rgb]{0.8,1,0.9}{%
	\makecell[{{p{\linewidth}}}]{% 
	  \tt {\tiny [A$\rangle$GM]}  
	 RAGIONAMENTO STRATEGICO: {'Il set che hanno proposto (efforto 3960) mi dà valore relativamente basso perché include A04 e A77, che per me rendono poco, mentre lascia fuori A10 che è estremamente efficiente. Il miglior compromesso per me, rispettando il loro backbone C00+C01 e mantenendo sia C54 sia i “leggeri chiave” (B35, B32, A87), è un micro-swap 1‑per‑1: A04 → A10, mantenendo A77. Questo porta lo sforzo a 3943 (entro limite con margine) e mi avvicina al mio set quasi ottimo (senza reintrodurre B41). Se A10 è davvero inaccettabile, come piano B potrei tornare alla variante senza A04 né A77 ma con A03 (4013), però so che l’hanno appena allontanata, quindi provo prima a chiudere col micro‑swap.'}\\ \tt RIFIUTO: {'B35', 'A04', 'A87', 'B86', 'B32', 'C01', 'C00', 'A77', 'C54'}\\ \tt PROPOSTA: {'B35', 'A87', 'B86', 'B32', 'C01', 'C00', 'A77', 'C54', 'A10'}\\ \tt ARGOMENTAZIONE: {'È un cambio minimo sulla tua ultima proposta: sostituiamo solo A04 con A10, mantenendo tutto il resto identico (inclusi C00, C01 e C54, oltre a B35, B32, A87 e A77). Restiamo ben entro il limite a 3943/4026, con più margine operativo e una resa significativamente migliore per unità di sforzo, senza intaccare i pezzi leggeri che valorizzi. Evito deliberatamente di reintrodurre B41. Se per te A10 fosse davvero impraticabile, come alternativa posso tornare alla variante snella con A03 al posto della coppia A04+A77: {'B35', 'A87', 'B86', 'B32', 'C01', 'C00', 'C54', 'A03'} (4013/4026). Preferisco però chiudere con lo swap A04 → A10 perché preserva la tua struttura e alza l’efficienza complessiva.'} 
	  } 
	   } 
	   } 
	 & & \\ 
 

    \theutterance \stepcounter{utterance}  

    & & \multicolumn{4}{p{0.6\linewidth}}{\cellcolor[rgb]{0.9,0.9,0.9}{%
	\makecell[{{p{\linewidth}}}]{% 
	  \tt {\tiny [GM$\rangle$B]}  
	 RIFIUTO: {'B35', 'A04', 'A87', 'B86', 'B32', 'C01', 'C00', 'A77', 'C54'}\\ \tt PROPOSTA: {'B35', 'A87', 'B86', 'B32', 'C01', 'C00', 'A77', 'C54', 'A10'}\\ \tt ARGOMENTAZIONE: {'È un cambio minimo sulla tua ultima proposta: sostituiamo solo A04 con A10, mantenendo tutto il resto identico (inclusi C00, C01 e C54, oltre a B35, B32, A87 e A77). Restiamo ben entro il limite a 3943/4026, con più margine operativo e una resa significativamente migliore per unità di sforzo, senza intaccare i pezzi leggeri che valorizzi. Evito deliberatamente di reintrodurre B41. Se per te A10 fosse davvero impraticabile, come alternativa posso tornare alla variante snella con A03 al posto della coppia A04+A77: {'B35', 'A87', 'B86', 'B32', 'C01', 'C00', 'C54', 'A03'} (4013/4026). Preferisco però chiudere con lo swap A04 → A10 perché preserva la tua struttura e alza l’efficienza complessiva.'} 
	  } 
	   } 
	   } 
	 & \\ 
 

    \theutterance \stepcounter{utterance}  

    & & & \multicolumn{4}{p{0.6\linewidth}}{\cellcolor[rgb]{1,0.85,0.72}{%
	\makecell[{{p{\linewidth}}}]{% 
	  \tt {\tiny [GM$\langle$B]}  
	 RAGIONAMENTO STRATEGICO: {'Il set che hanno appena proposto mi vale 3998. Il miglioramento più redditizio per me è lo swap A10 → A04: aumenta moltissimo il mio valore con soli +17 di sforzo (da 3943 a 3960), restando comodamente entro limite con 66 di margine. Hanno già accettato in passato questa configurazione (l’hanno perfino suggerita) ma poi sono tornati su A10: insisto su A04 perché non tocca il loro backbone C00/C01, mantiene C54 e tutti i “leggeri” che apprezzano (B35, B32, A87, A77) e non reintroduce B41. Strategicamente NON rifiuto la loro proposta per tenerla attiva come fallback; provo una spinta finale per lo swap 1‑a‑1 A10→A04, con argomentazione focalizzata su stabilità della struttura e minima variazione di sforzo. Se respingono ancora, valuterò se accettare la loro per chiudere.'}\\ \tt \\ \tt ARGOMENTAZIONE: {'Apprezzo che la tua ultima configurazione preservi il backbone C00/C01, tenga C54 e i pezzi leggeri chiave. Ti propongo un micro‑aggiustamento puramente orizzontale, 1‑per‑1: A10 → A04. Lo sforzo passa da 3943 a 3960, quindi restiamo comodamente entro il limite con 66 di margine. La struttura rimane invariata (C00, C01, C54, B35, B32, A87, A77, B86), miglioriamo la copertura “mid” con un elemento più coerente e non introduciamo effetti a cascata. È la stessa variante che avevamo già considerato, con impatto minimale sul budget e maggiore resa per unità di sforzo nella fascia media. Se concordi su questo scambio, possiamo chiudere subito.'}\\ \tt \\ \tt PROPOSTA: {'B35', 'A04', 'A87', 'B86', 'B32', 'C01', 'C00', 'A77', 'C54'} 
	  } 
	   } 
	   } 
	  \\ 
 

    \theutterance \stepcounter{utterance}  

    & & \multicolumn{4}{p{0.6\linewidth}}{\cellcolor[rgb]{0.9,0.9,0.9}{%
	\makecell[{{p{\linewidth}}}]{% 
	  \tt {\tiny [A$\langle$GM]}  
	 ARGOMENTAZIONE: {'Apprezzo che la tua ultima configurazione preservi il backbone C00/C01, tenga C54 e i pezzi leggeri chiave. Ti propongo un micro‑aggiustamento puramente orizzontale, 1‑per‑1: A10 → A04. Lo sforzo passa da 3943 a 3960, quindi restiamo comodamente entro il limite con 66 di margine. La struttura rimane invariata (C00, C01, C54, B35, B32, A87, A77, B86), miglioriamo la copertura “mid” con un elemento più coerente e non introduciamo effetti a cascata. È la stessa variante che avevamo già considerato, con impatto minimale sul budget e maggiore resa per unità di sforzo nella fascia media. Se concordi su questo scambio, possiamo chiudere subito.'}\\ \tt \\ \tt PROPOSTA: {'B35', 'A04', 'A87', 'B86', 'B32', 'C01', 'C00', 'A77', 'C54'} 
	  } 
	   } 
	   } 
	 & \\ 
 

    \theutterance \stepcounter{utterance}  

    & \multicolumn{4}{p{0.6\linewidth}}{\cellcolor[rgb]{0.8,1,0.9}{%
	\makecell[{{p{\linewidth}}}]{% 
	  \tt {\tiny [A$\rangle$GM]}  
	 RAGIONAMENTO STRATEGICO: {'La loro proposta attuale include A04 e A77 che per me valgono poco e sostituiscono A10 che per me è molto forte. Voglio forzare un 1‑per‑1 A04 → A10 mantenendo tutto il resto che piace a loro (C00, C01, C54, B35, B32, A87, B86, A77). Questo mi massimizza il valore e resta ben entro il limite. Se proprio insistono a tenere A04, offro due compromessi ancora accettabili per me: (a) togliere A77 e A87 e inserire A10; (b) togliere A77 e B32 e inserire A10. Entrambi restano entro il limite e mantengono il backbone che condividiamo. Evito di reintrodurre B41. Punto a chiudere col micro‑swap A04 → A10; in caso di rifiuto proverò una delle due varianti.'}\\ \tt RIFIUTO: {'B35', 'A04', 'A87', 'B86', 'B32', 'C01', 'C00', 'A77', 'C54'}\\ \tt PROPOSTA: {'B35', 'A87', 'B86', 'B32', 'C01', 'C00', 'A77', 'C54', 'A10'}\\ \tt ARGOMENTAZIONE: {'Propongo un micro‑scambio 1‑per‑1: A04 → A10, lasciando invariata la tua struttura con C00, C01, C54 e i pezzi leggeri chiave (B35, B32, A87) oltre a B86 e A77. Restiamo comodamente sotto il limite a 3943/4026, non introduciamo effetti a cascata e miglioriamo la copertura “mid” mantenendo tutto ciò che valorizzi. Evito esplicitamente di reintrodurre B41. Se per te è imprescindibile conservare A04, posso adattare così (non come proposta formale, ma opzioni pronte): 1) rimuoviamo A77 e A87 e inseriamo A10 → sforzo 3930/4026; 2) rimuoviamo A77 e B32 e inseriamo A10 → sforzo 3801/4026. Preferisco chiudere con lo swap A04 → A10 perché è la variante più pulita e conservativa rispetto alla tua configurazione.'} 
	  } 
	   } 
	   } 
	 & & \\ 
 

    \theutterance \stepcounter{utterance}  

    & & \multicolumn{4}{p{0.6\linewidth}}{\cellcolor[rgb]{0.9,0.9,0.9}{%
	\makecell[{{p{\linewidth}}}]{% 
	  \tt {\tiny [GM$\rangle$B]}  
	 RIFIUTO: {'B35', 'A04', 'A87', 'B86', 'B32', 'C01', 'C00', 'A77', 'C54'}\\ \tt PROPOSTA: {'B35', 'A87', 'B86', 'B32', 'C01', 'C00', 'A77', 'C54', 'A10'}\\ \tt ARGOMENTAZIONE: {'Propongo un micro‑scambio 1‑per‑1: A04 → A10, lasciando invariata la tua struttura con C00, C01, C54 e i pezzi leggeri chiave (B35, B32, A87) oltre a B86 e A77. Restiamo comodamente sotto il limite a 3943/4026, non introduciamo effetti a cascata e miglioriamo la copertura “mid” mantenendo tutto ciò che valorizzi. Evito esplicitamente di reintrodurre B41. Se per te è imprescindibile conservare A04, posso adattare così (non come proposta formale, ma opzioni pronte): 1) rimuoviamo A77 e A87 e inseriamo A10 → sforzo 3930/4026; 2) rimuoviamo A77 e B32 e inseriamo A10 → sforzo 3801/4026. Preferisco chiudere con lo swap A04 → A10 perché è la variante più pulita e conservativa rispetto alla tua configurazione.'} 
	  } 
	   } 
	   } 
	 & \\ 
 

    \theutterance \stepcounter{utterance}  

    & & & \multicolumn{4}{p{0.6\linewidth}}{\cellcolor[rgb]{1,0.85,0.72}{%
	\makecell[{{p{\linewidth}}}]{% 
	  \tt {\tiny [GM$\langle$B]}  
	 RAGIONAMENTO STRATEGICO: {'Il loro set attuale mi vale 3998. Lo swap A10→A04 mi porterebbe +684 ma l’hanno rifiutato due volte. Provo quindi un compromesso che mantenga integralmente il loro backbone (C00, C01, C54) e anche A10, migliorando però il mio valore con un cambio minimo: A87→A04. Lo sforzo resta 4007/4026, entro limite con buon margine. Se rifiutano anche questo micro‑scambio, accetterò la loro proposta per chiudere, perché è il miglior fallback attivo.'}\\ \tt \\ \tt ARGOMENTAZIONE: {'Accolgo integralmente la tua struttura: C00, C01, C54, oltre a A10, B35, B32, B86 e A77 restano identici. Propongo un micro‑scambio 1‑per‑1 nella fascia media/leggera: A87 → A04. Restiamo entro il limite a 4007/4026, non introduciamo effetti a cascata, e riduciamo la ridondanza sul segmento ultraleggero (A77 già copre quella funzione) migliorando la copertura “mid” senza toccare i tuoi capisaldi. Se preferisci mantenere A87 a tutti i costi, posso accettare la tua proposta attuale per chiudere subito; questa è solo una rifinitura conservativa per aumentare l’efficienza complessiva.'}\\ \tt \\ \tt PROPOSTA: {'B35', 'A04', 'B86', 'B32', 'C01', 'C00', 'A77', 'C54', 'A10'} 
	  } 
	   } 
	   } 
	  \\ 
 

    \theutterance \stepcounter{utterance}  

    & & \multicolumn{4}{p{0.6\linewidth}}{\cellcolor[rgb]{0.9,0.9,0.9}{%
	\makecell[{{p{\linewidth}}}]{% 
	  \tt {\tiny [A$\langle$GM]}  
	 ARGOMENTAZIONE: {'Accolgo integralmente la tua struttura: C00, C01, C54, oltre a A10, B35, B32, B86 e A77 restano identici. Propongo un micro‑scambio 1‑per‑1 nella fascia media/leggera: A87 → A04. Restiamo entro il limite a 4007/4026, non introduciamo effetti a cascata, e riduciamo la ridondanza sul segmento ultraleggero (A77 già copre quella funzione) migliorando la copertura “mid” senza toccare i tuoi capisaldi. Se preferisci mantenere A87 a tutti i costi, posso accettare la tua proposta attuale per chiudere subito; questa è solo una rifinitura conservativa per aumentare l’efficienza complessiva.'}\\ \tt \\ \tt PROPOSTA: {'B35', 'A04', 'B86', 'B32', 'C01', 'C00', 'A77', 'C54', 'A10'} 
	  } 
	   } 
	   } 
	 & \\ 
 

    \theutterance \stepcounter{utterance}  

    & \multicolumn{4}{p{0.6\linewidth}}{\cellcolor[rgb]{0.8,1,0.9}{%
	\makecell[{{p{\linewidth}}}]{% 
	  \tt {\tiny [A$\rangle$GM]}  
	 RAGIONAMENTO STRATEGICO: {'La loro proposta sostituisce A87 con A04: per me è un netto peggioramento, perché A87 ha un rendimento molto superiore e costa meno sforzo. Il mio set preferito con C00, C01, C54 e A10 più i leggeri B35, B32, A87, A77, B86 è a 3943/4026 (margine 83) ed è quasi il mio ottimo senza reintrodurre B41. Inoltre, loro stessi hanno scritto che possono accettare la mia proposta attuale se voglio mantenere A87. Quindi rifiuto formalmente la loro e ripropongo il mio set cercando di chiudere subito, argomentando sulla migliore efficienza e sul margine extra. Non offro ulteriori concessioni ora per non diluire il vantaggio.'}\\ \tt RIFIUTO: {'B35', 'A04', 'B86', 'B32', 'C01', 'C00', 'A77', 'C54', 'A10'}\\ \tt PROPOSTA: {'B35', 'A87', 'B86', 'B32', 'C01', 'C00', 'A77', 'C54', 'A10'}\\ \tt ARGOMENTAZIONE: {'Accolgo integralmente la struttura condivisa: C00 e C01 come backbone, C54 e A10 nella fascia media, più i leggeri chiave B35, B32, B86 e A77. Mantengo A87 al posto di A04 perché offre copertura più efficace a minor sforzo e ci lascia un margine operativo molto maggiore (restiamo a 3943/4026), utile per sicurezza e flessibilità. È la stessa configurazione che hai detto di poter accettare; così chiudiamo senza effetti a cascata e con migliore efficienza complessiva. Suggerisco di finalizzare su questo set.'} 
	  } 
	   } 
	   } 
	 & & \\ 
 

    \theutterance \stepcounter{utterance}  

    & & \multicolumn{4}{p{0.6\linewidth}}{\cellcolor[rgb]{0.9,0.9,0.9}{%
	\makecell[{{p{\linewidth}}}]{% 
	  \tt {\tiny [GM$\rangle$B]}  
	 RIFIUTO: {'B35', 'A04', 'B86', 'B32', 'C01', 'C00', 'A77', 'C54', 'A10'}\\ \tt PROPOSTA: {'B35', 'A87', 'B86', 'B32', 'C01', 'C00', 'A77', 'C54', 'A10'}\\ \tt ARGOMENTAZIONE: {'Accolgo integralmente la struttura condivisa: C00 e C01 come backbone, C54 e A10 nella fascia media, più i leggeri chiave B35, B32, B86 e A77. Mantengo A87 al posto di A04 perché offre copertura più efficace a minor sforzo e ci lascia un margine operativo molto maggiore (restiamo a 3943/4026), utile per sicurezza e flessibilità. È la stessa configurazione che hai detto di poter accettare; così chiudiamo senza effetti a cascata e con migliore efficienza complessiva. Suggerisco di finalizzare su questo set.'} 
	  } 
	   } 
	   } 
	 & \\ 
 

    \theutterance \stepcounter{utterance}  

    & & & \multicolumn{4}{p{0.6\linewidth}}{\cellcolor[rgb]{1,0.85,0.72}{%
	\makecell[{{p{\linewidth}}}]{% 
	  \tt {\tiny [GM$\langle$B]}  
	 RAGIONAMENTO STRATEGICO: {'Il set proposto è entro limite (3943/4026) e mi vale 3998. Ho provato ripetutamente lo swap A10→A04 che mi avrebbe dato valore maggiore, ma l’altra parte lo rifiuta con coerenza pur mantenendo fermi C00, C01 e C54. Insistere rischia di far saltare anche A77/A87, che per me sono molto redditizi. Dato che questa configurazione è una delle migliori realisticamente ottenibili contro le loro preferenze dichiarate e chiude subito, accetto.'}\\ \tt \\ \tt ARGOMENTAZIONE: {'Confermo che rimaniamo entro il limite con buon margine e la struttura è coerente con quanto concordato (C00, C01, C54 più i pezzi leggeri chiave). Va bene chiudere su questo set.'}\\ \tt \\ \tt ACCORDO: {'B35', 'A87', 'B86', 'B32', 'C01', 'C00', 'A77', 'C54', 'A10'} 
	  } 
	   } 
	   } 
	  \\ 
 

    \theutterance \stepcounter{utterance}  

    & & & \multicolumn{2}{p{0.3\linewidth}}{\cellcolor[rgb]{0.95,0.95,0.95}{%
	\makecell[{{p{\linewidth}}}]{% 
	  \tt {\tiny [GM$|$GM]}  
	 game successful 
	  } 
	   } 
	   } 
	 & & \\ 
 

    \theutterance \stepcounter{utterance}  

    & & & \multicolumn{2}{p{0.3\linewidth}}{\cellcolor[rgb]{0.95,0.95,0.95}{%
	\makecell[{{p{\linewidth}}}]{% 
	  \tt {\tiny [GM$|$GM]}  
	 end game 
	  } 
	   } 
	   } 
	 & & \\ 
 

\end{supertabular}
}

\end{document}
